% concatenated from Arxiv.tex
\documentclass[12pt]{article}

\usepackage{amsmath,amsthm}
\usepackage{mathtools}
\usepackage{xcolor}
\usepackage{amssymb}
\usepackage{amsfonts}
\usepackage{graphicx}
\usepackage{booktabs}
\usepackage{units}
\usepackage{xcolor}
\usepackage{comment}
\usepackage{caption}
\usepackage{subcaption} % for side-by-side images
\usepackage{dblfloatfix}   % patches LaTeX’s double-column float rules
\usepackage{algorithm}
\usepackage{algorithmicx}
\usepackage{algpseudocode}
\algrenewcommand\algorithmicindent{0.5em}
\usepackage{commath}

\usepackage{hyperref}
\hypersetup{
    colorlinks=true,
    linkcolor=blue,     
    urlcolor=blue,
}
\usepackage{xparse}

\usepackage{tikz}
\usepackage{pgfplots}
\usepackage{appendix}
\let\labelindent\relax
\usepackage{enumitem}

\newcommand{\R}{\mathbb{R}}
\newcommand{\N}{\mathbb{N}}
\newcommand{\B}{\mathbb{B}}
\newcommand{\mbb}[1]{\mathbb{#1}}
\newcommand{\mrm}[1]{\operatorname{#1}}
\newcommand{\mc}[1]{\mathcal{#1}}
\newcommand{\diag}[1]{\mrm{diag}\big({#1}\big)}
\floatname{algorithm}{Pseudo code}
\DeclareMathOperator{\ran}{ran}
\DeclareMathOperator{\sym}{sym}

\newtheorem{definition}{Definition}

\allowdisplaybreaks
% Define a shared counter for all theorem-like environments:
\newcounter{thmcounter}[section]  % Resets with each new section
\renewcommand{\thethmcounter}{\thesection.\arabic{thmcounter}}

% Helper macro to define an environment with a QED symbol and an optional title.
% #1 = name of the environment (e.g., theorem, lemma)
% #2 = printed label (e.g., Theorem, Lemma)
\newcommand{\defthmwithqed}[2]{%
  \NewDocumentEnvironment{#1}{ o }{%
    \refstepcounter{thmcounter}% increment counter and allow referencing
    \begin{trivlist}%
      \item[\hskip \labelsep \bfseries #2~\thethmcounter%
        \IfValueT{##1}{\ (##1)}.]%
  }{%
    \hfill $\square$%
    \end{trivlist}%
  }%
}


% Now define your environments:
\newtheorem{theorem}{Theorem}
\newtheorem{lemma}{Lemma}
\newtheorem{assumption}{Assumption}
\newtheorem{proposition}{Proposition}
\newtheorem{prob}{Problem}
\newtheorem{corollary}{Corollary}
\newtheorem{remark}{Remark}
\newtheorem{example}{Example}





\usepackage{tikz}
\usetikzlibrary{arrows.meta,positioning,calc,decorations.pathmorphing}

\newenvironment{smallpmatrix}%          environment name
{\left(\begin{smallmatrix}}%            begin code
{\end{smallmatrix}\right)}%             end code

\newenvironment{smallbmatrix}%          environment name
{\left[\begin{smallmatrix}}%            begin code
{\end{smallmatrix}\right]}%       


\title{Dissipativity properties of a class of nonlinear time-delay systems via Bessel--Legendre inequalities}
\author{Ikram El Haskouki\footnote{AIMCE Laboratory, ENSAM, Hassan II University of Casablanca, 150 Nile Boulevard, Casablanca 20670, Morocco and University of Wuppertal, Gaußstraße 20, 42119 Wuppertal, Germany (e-mail: \texttt{ikram.elhaskouki[at]gmail.com})} and Hannes Gernandt\footnote{University of Wuppertal, Gaußstraße 20, 42119 Wuppertal and Fraunhofer IEG, Fraunhofer Research Institution for Energy Infrastructures and Geotechnologies IEG, 03046 Cottbus, Germany  (e-mail: \texttt{gernandt[at]uni-wuppertal.de})} 
}


\date{\today}

\begin{document}

\maketitle
%\tableofcontents

\begin{abstract}
Time delays are inherent in many physical and engineered systems and can significantly affect 
their stability and performance. In this work, we investigate the dissipativity of a class of nonlinear time-delay systems with multiple discrete delays and derive sufficient conditions for both delay-dependent and delay-independent dissipativity using Bessel--Legendre inequalities. For linear systems, the resulting  dissipativity conditions are expressed in terms of linear matrix inequalities (LMIs) which can be solved numerically to obtain Lyapunov--Krasovskii-type storage functions. 
\end{abstract}


\maketitle


\section{Introduction}
\label{sec:introduction}
Time delays arise naturally in many physical, biological, and engineering systems, for instance due to transport phenomena, communication networks, measurement processes, or delayed feedback. Their presence may significantly influence both stability and performance and can lead to qualitatively different behavior compared with delay-free systems. For this reason, the analysis and control of time-delay systems has been an active research area for several decades; see, for instance, \cite{Fridman2014,Scholl2024} and the references therein.

Dissipativity provides a powerful framework for the analysis and design of dynamical systems through input--output energy balance relations. It was originally introduced by Willems \cite{willems1972dissipative} and further developed by Hill and Moylan \cite{hill1976stability}. In this framework, the change of a suitable storage function is bounded by a prescribed supply rate, which describes the energy exchanged between the system and its environment. Important system properties such as impedance passivity and scattering passivity can be formulated as special cases of dissipativity by choosing appropriate quadratic supply rates. This makes dissipativity particularly useful for stability analysis, robust control design, and energy-based modeling frameworks, and in particular in form of port-Hamiltonian systems \cite{Brogliato2007,BreitHU24,haskouki2026delaydependentpassivitystabilitylinear}.

In the presence of time delays, dissipativity analysis becomes more involved, since the system state depends not only on its current value but also on its past history. Consequently, classical Lyapunov functions are generally not sufficient to capture the influence of the delayed state. Instead, Lyapunov--Krasovskii storage functionals are commonly employed, as they explicitly incorporate information on the state history \cite{Fridman2014}. A central difficulty in this approach is the estimation of integral terms involving delayed states and their derivatives. Jensen's inequality has been widely used for this purpose, but it may lead to conservative conditions. To reduce this conservatism, several refined integral inequalities have been proposed, including Wirtinger-based inequalities \cite{seuret2013jensen,seuret2013wirtinger}.

Recently, Bessel--Legendre inequalities have emerged as a powerful extension of Jensen and Wirtinger-type estimates \cite{seuret2015hierarchy}. By exploiting orthogonal projections onto Legendre polynomial bases, these inequalities provide a hierarchy of increasingly accurate estimates for integral terms while preserving a tractable structure. They are therefore well suited for deriving less conservative Lyapunov--Krasovskii conditions for time-delay systems.

Motivated by these developments, this paper investigates dissipativity properties of a class of nonlinear time-delay systems with multiple constant delays. We focus on general quadratic supply rates, which include impedance and scattering passivity as important special cases. The proposed approach combines Lyapunov--Krasovskii storage functionals with Bessel--Legendre inequalities in order to derive sufficient delay-dependent dissipativity conditions. A delay-independent criterion is obtained as a special case by considering a simplified storage functional. Moreover, for linear time-delay systems, the resulting conditions are expressed as linear matrix inequalities, which enables efficient numerical verification. Recent results on delay-dependent stability and passivity of linear port-Hamiltonian delay systems \cite{haskouki2026delaydependentpassivitystabilitylinear} further motivate the study of such dissipativity conditions in the presence of delays.

The main contributions of this paper are summarized as follows.
\begin{itemize}
\item We derive sufficient delay-dependent dissipativity conditions for nonlinear systems with multiple constant delays by means of Lyapunov--Krasovskii storage functionals and Bessel--Legendre inequalities.
\item We obtain delay-independent dissipativity conditions as a special case based on a simplified storage functional.
\item We specialize the proposed framework to linear time-delay systems and derive tractable LMI conditions for dissipativity with respect to general quadratic supply rates.
\end{itemize}
The paper is organized as follows. After introducing the notation, Section~\ref{sec:dissipative} recalls the required properties of Legendre polynomials and the Bessel--Legendre inequality. It then establishes delay-dependent and delay-\linebreak independent dissipativity conditions for nonlinear time-delay systems. Section~\ref{sec:linear} treats the linear case and derives corresponding LMI conditions. Finally, the paper concludes with a summary and an outlook on future research directions in Section~\ref{sec:conclusion}.

\section*{Notation}
Throughout the paper, $I_n$ stands for the $n \times n$ identity matrix, and $0$ denotes a zero matrix of appropriate dimension.
For $x \in  C([-\tau,0],\mathbb{R}^{n}),$ and 
$t \geqslant 0,$ one defines the function  
$x_t \in  C([-\tau,0],\mathbb{R}^{n})$ by 
$x_t(\theta) = x(t+\theta),$ for all $\theta 
\in [-\tau, 0].$ The space $C([-
\tau,0],\mathbb{R}^{n})$ of continuous functions defined from $[-
\tau,0]$ to $\mathbb{R}^{n}$  is equipped with the norm
$ \Vert \psi \Vert_{\infty} = \sup\limits_{-\tau \leqslant 
s \leqslant 0} \Vert \psi(s)\Vert,$ where  $\|\cdot\|$ denotes the Euclidean norm in $\mathbb{R}^n.$ The notation $P>0$ ($P \ge 0$), for $P \in \mathbb{R}^{n \times 
n}$ means that $P$ is symmetric and positive definite (positive 
semidefinite) and $\mathrm{sym}(A) := \frac{1}{2}(A + A^\top)$ denotes the symmetric part of a matrix $A \in \mathbb{R}^{n \times n }.$
\section{Dissipativity of nonlinear time-delay systems}
\label{sec:dissipative}
In this section, we consider a class of nonlinear time-delay systems that are given by 
\begin{align}
\nonumber 
\dot x(t) &= f(x(t),x(t-\tau_1),\ldots, x(t-\tau_q)) 
  + g(x(t))u(t),\quad t\geq 0,\\
x(t) &= \phi(t),\quad  t \in [-\tau_q , 0], \label{eq:cnt_time_delay_system_nonlinear}
\\
y(t) &= h(x(t)),\quad t \geq  0,\nonumber 
\end{align}
where $u : \mathbb{R} \to \mathbb{R}^m,$ $y : \mathbb{R} \to \mathbb{R}^m,$ $x : \mathbb{R} \to \mathbb{R}^n $ are the input, output, and state of the system, $f:\mathbb{R}^{n(q+1)}\rightarrow\mathbb{R}^n$ and $g:\mathbb{R}^n\rightarrow\mathbb{R}^{n\times m}$ all locally Lipschitz continuous and $\tau_q \geqslant\ldots \geqslant \tau_1 > 0$ for some natural number $q\geq 1$.


\begin{definition}
A delay system \eqref{eq:cnt_time_delay_system_nonlinear} with delays $\tau_q \geqslant\ldots \geqslant \tau_1 > 0$ is called \emph{dissipative} w.r.t. the continuous supply rate $s:\mathbb{R}^m\times\mathbb{R}^m\rightarrow\mathbb{R}$ if there exists a storage function $\mathcal{S}:C^1([-\tau_q,0],\mathbb{R}^n)\rightarrow[0,\infty)$ with $\mathcal{S}(0)=0$ such that the following dissipation inequality holds
\begin{align}
    \label{eq:dissip_ineq_cont_time}
\mathcal{S}(x_{t_1})-\mathcal{S}(x_{t_0})\leq \int_{t_0}^{t_1}s(y(s),u(s))\mathrm{d}s,\quad \text{for all}~ 0\leq t_0\leq t_1<\infty,
\end{align}
\end{definition}
and all initial histories {\footnotesize $\phi \in C^1([-\tau_q,0],\mathbb{R}^n)$}, where {\footnotesize
$x_t=x_{[t-\tau_q,t]}\in C([-\tau_q,0],\mathbb{R}^n)$
} is given by $x_t(\theta)=x(t+\theta)$, $\theta\in[-\tau_q,0]$, where $x$ is the solution to \eqref{eq:cnt_time_delay_system_nonlinear} for some history $\phi$. 

If the dissipation inequality \eqref{eq:dissip_ineq_cont_time} holds for all $\tau\geq 0$, then we call the system \eqref{eq:cnt_time_delay_system_nonlinear} \emph{dissipative independent of the delay.}

In this paper, we focus on quadratic supply rates of the form 
\[s(y,u) =\label{eq:supply_rate}
\begin{bmatrix} y\\ u \end{bmatrix}^\top
\begin{bmatrix} S_{11} & S_{12} \\ S_{12}^\top & S_{22} \end{bmatrix}
\begin{bmatrix} y\\ u \end{bmatrix},\]
\[S_{11}=S_{11}^\top\in\mathbb{R}^{m \times m}, \quad S_{12}\in\mathbb{R}^{m\times m},\quad S_{22}=S_{22}^\top \in \mathbb{R}^{m \times m}.
\]
\vspace{1mm}
\\
Systems which are dissipative with respect to the supply rate 
\[
s_{imp}(y,u)=\frac12\begin{bmatrix}
    y\\ u
\end{bmatrix}^\top\begin{bmatrix}
    0&I_m\\ I_m&0
\end{bmatrix}\begin{bmatrix}
    y\\ u
\end{bmatrix},\quad s_{sca}(y,u)=\frac12\begin{bmatrix}
    y\\ u
\end{bmatrix}^\top\begin{bmatrix}
    -I_m &0\\ 0&I_m
\end{bmatrix}\begin{bmatrix}
    y\\ u
\end{bmatrix}
\]
are called \emph{impedance} and \emph{scattering passive} respectively.
\vspace{1mm}\\
\subsection{Legendre Polynomials and the Bessel--Legendre inequality}
In the following, we provide a brief overview of the Legendre polynomials and their relevant properties.
\begin{definition}
The Legendre polynomials considered over the interval $[-\tau_i, 0]$ are defined by: 
\begin{equation}
    \forall k \in \mathbb{N},\,i=1,\ldots,q: \quad L_k^i(s) = (-1)^k \sum_{l=0}^k (-1)^l \binom{k}{l} \binom{k+l}{l} \left( \frac{s+\tau_i}{\tau_i} \right)^l.
\end{equation}
\end{definition}


The Legendre polynomials have the following properties, see e.g. \cite{seuret2015hierarchy}.

\begin{lemma}\rm\ \label{Lemma:Legendre_polynomials_properties}
Consider the Legendre polynomials $L_k^i$ for $i=1,\ldots,q$ and $k \in \mathbb{N}$. Then the following holds 
\begin{itemize}
\item \textit{Orthogonality:} For all $(k, l) \in \mathbb{N}^2$ and all $i=1,\ldots,q$,
\begin{equation}\nonumber
    \int_{-\tau_i}^0 L_k^i(s)L_l^i(s) \, \mathrm{d}s = 
    \begin{cases} 
        0, & k \neq l, \\
        \frac{\tau_i}{2k+1}, & k = l.
    \end{cases}
\end{equation}
\item \textit{Boundary conditions:}
$    \forall k \in \mathbb{N}, \quad L_k^i(0) = 1, \quad L_k^i(-\tau_i) = (-1)^k$,\\ $i=1,\ldots,q$.

\item \textit{Differentiation:} The derivative $\dot{L}^i_k$ is given by:
\begin{equation}\nonumber
    \dot{L}^i_k(s) = 
    \begin{cases} 
        0, & k = 0, \\
        \displaystyle \sum_{j=0}^{k-1} \frac{2j+1}{\tau_i} \left(1 - (-1)^{k+j}\right) L^i_j(s), & k \ge 1.
    \end{cases}
\end{equation}
\end{itemize}
\end{lemma}
The following lemma which is only a minor modification of \cite[Corollary 4]{seuret2015hierarchy}, known as the Bessel–Legendre inequality, is 
established using Legendre polynomials. 
\begin{lemma}\label{lemma:legendre}
Let $x$ be such that 
$x \in C^1([-\tau_i,0],\mathbb{R}^{n})$, $R_i\geqslant0$ and $\tau_i> 0$. 
Then, the 
integral inequality
\begin{equation}\label{eq:bessel}
\int_{-\tau_i}^{0} \dot{x}(u)^\top R_i\dot{x}
(u)du \geq \frac{1}{\tau_i}\xi_{N}^i(t)^{\top} \left[ 
\sum_{k=0}^{N} (2k+1) \Gamma_{N}(k)^\top R_i 
\Gamma_N(k) \right] \xi^i_{N}(t), 
\end{equation}
holds, for all integers $N \in 
\mathbb{N}$ with $\Omega^i_k(t)= \int_{-\tau_i}^0L_k^i(s)x_t(s)\mathrm{d}s, \hspace{2mm} k\in \mathbb{N}$ and 
\begin{align*}
\xi_{N}^i(t) &= 
\begin{cases} 
[x_t(0)^\top \ x_t(-\tau_i)^\top]^\top, & \text{if } N = 
0, \\
\left[ x_t(0)^\top \quad x_t(-\tau_i)^\top \quad 
\frac{1}{\tau_i}\Omega_0^{i}(t)^\top\quad \dots 
\quad \frac{1}{\tau_i}\Omega_{N-1}^{i}(t)^\top
\right]^\top, & \text{if } N \geq 1,
\end{cases}\\
\Gamma_N(k) &= 
\begin{cases} 
\begin{bmatrix} I & -I \end{bmatrix}, 
& \text{if } N = 0, \\
\begin{bmatrix} I & (-1)^{k+1}I & 
\gamma_{Nk}^0 I & \dots & 
\gamma_{Nk}^{N-1}I \end{bmatrix}, & 
\text{if } N \geq  1,
\end{cases}\\
\gamma_{Nk}^j &= 
\begin{cases} 
-(2j+1)(1 - (-1)^{k+j}), & \text{if } 
j \leq k, \\
0, & \text{if } j > k.
\end{cases}
\end{align*}
\end{lemma}
\begin{proof}
The proof is identical to the one given 
in \cite{seuret2015hierarchy}. The assumption that 
$R_i$ is positive definite can be relaxed to positive semidefinite 
without affecting the result by following the same argument as in the proof of  \cite[Lemma 2.2]{haskouki2026delaydependentpassivitystabilitylinear}. 
\end{proof}
\begin{remark}
The Bessel–Legendre inequality is a generalization of classical Jensen and Wirtinger-based
inequalities, which are obtained for $N=0$ and $N=1$, respectively, which were used for delay systems in \cite{seuret2013wirtinger,Fridman2014}.
\end{remark}
\subsection{Dissipativity via nonlinear matrix inequalities}
In the following, we aim to show delay-dependent dissipativity with respect to the quadratic supply rate  \eqref{eq:supply_rate} by using a modified Lyapunov--Krasovskii function from \cite{seuret2015hierarchy} as a storage function on $C^1([-\tau_q,0],\mathbb{R}^n).$ 

We consider the following storage function :
\begin{equation}
\begin{split}\label{eq:storage_fixed_delay_nonlinear}
\mathcal{S}_N(x_t) &= \mathcal{S}_1(x(t))+\hat{x}_{N}(t)^\top P_{N} \hat{x}_{N}(t) + 
\sum_{i=1}^q \int_{t-\tau_i}^{t} \Theta_i( x(s) )\mathrm{d}s\\
&+ \sum_{i=1}^q \tau_i \int_{t-\tau_i}^{t} \int_{\theta}^{t} \dot{x}(s)^\top R_i\dot{x}(s) 
\mathrm{d}s \mathrm{d}\theta,
\end{split}
\end{equation}
with the following extra-states 
\begin{equation*}
\hat{x}_{N}(t) = \begin{bmatrix}
\Omega_0^{1}(t)^\top & 
\ldots &\Omega_{N-1}^{1}(t)^\top &
\ldots &
\Omega_{0}^{q}(t)^\top & \ldots &  \Omega_{N-1}^{q}(t)^\top
\end{bmatrix}^\top, \hspace{2mm} \text{if} \hspace{2mm} N \geq 1,
\end{equation*} 
and $\hat{x}_0(t)=0$ if $N=0,$ 
for some continuously differentiable $\mathcal{S}_1:\mathbb{R}^n\rightarrow[0,\infty)$, $P_N\in \mathbb{R}^{qnN \times qnN}$ positive semi-definite, continuous positive semi-definite function 
$\Theta_i:\mathbb{R}^n\rightarrow[0,\infty)$ and $R_i\geq 0,$ for $i \in \{1,\ldots,q\}.$ Furthermore, we consider the following matrix functions for $x_0,x_i\in\mathbb{R}^n$  and $\xi_N^i\in \mathbb{R}^{n(N+2)}$ 
 {\footnotesize \begin{align*}
 \Psi_{xx}(x_0, \ldots,x_q,\xi^1_{N},\ldots,\xi^q_N)&:=\nabla\mathcal{S}_1(x_0)^\top f(x_0,\ldots,x_q)+  2\zeta(t)^\top\mathrm{sym}(G^\top_NP_N H_N)\zeta(t)\\
 &+\sum_{i=1}^q\tau_i^2 f(x_0,x_1,\ldots,x_q)^{\top} R_if(x_0,x_1,\ldots,x_q)\\
 &+ \sum_{i=1}^q \Theta_i(x_0)-\sum_{i=1}^q\Theta_i(x_i)\\
&{-\zeta(t)^\top \mathbf{\Gamma_N}^\top \mathbf{R_N}\mathbf{\Gamma_N}\zeta(t)}\\
&-h(x_0)^\top S_{11}h(x_0).
\\ 
  \Psi_{xu}(x_0,\ldots,x_q)&:=\tfrac12 \nabla\mathcal{S}_1(x_0)^\top g(x_0)+\sum_{i=1}^q \tau_i ^2f(x_0,x_1,\ldots,x_q)^\top R_i g(x_0)\\
  &-h(x_0)^\top S_{12}\\
  \Psi_{uu}(x_0)&:=\sum_{i=1}^q \tau_i^2g(x_0)^\top R_i g(x_0)-S_{22},\\
\zeta(t)&=\left[\begin{smallmatrix}x_t(0)^\top \quad x_t(-\tau_1)^\top \ldots\,  x_t(-\tau_q)^\top\quad \left[\frac{1}{\tau_i}\Omega_0^{i}(t)^\top \ldots\,  \frac{1}{\tau_i}\Omega_{N-1}^{i}(t)^\top\right]_{i=1}^q 
\end{smallmatrix}\right]^\top,\\
G_N&=\begin{bmatrix} 
0&0&0&0\\
  0& 0 & \rm diag( \tau_i I_{nN})_{i=1}^q&0 \end{bmatrix},\hspace{2mm} H_N=\begin{bmatrix}
\tilde \Gamma_{N,1}(0)\\ \vdots\\
\tilde \Gamma_{N,1}(N-1) \\
\tilde \Gamma_{N,2}(0)\\
\vdots
\\\tilde \Gamma_{N,q}(N-1) 
\end{bmatrix},\\
R_{N,i}&=\rm diag(R_i,3R_i, \ldots, (2N+1)R_i),\\
 \mathbf{R_N}&=\rm diag(R_{N,1}, \ldots, R_{N,q}),\\
\mathbf{\Gamma_N} &= \begin{bmatrix} \Gamma_{N,1} \\  \Gamma_{N,2} \\ \vdots \\ \Gamma_{N,q} \end{bmatrix},  \hspace{2mm} \Gamma_{N,i} = \begin{bmatrix} \tilde\Gamma_{N,i}(0) \\ \tilde \Gamma_{N,i}(1) \\ \vdots \\ \tilde\Gamma_{N,i}(N) \end{bmatrix},\\
\tilde \Gamma_{0,i}(0)&=
\begin{bmatrix}
I & 0 & \ldots &0  &
\underbrace{-I}_{(i+1)\text{-th block}} &
0 & \ldots & 0
\end{bmatrix}\\
\text{and  if }  N\geq 1, \text{we have}
\end{align*}}
$\tilde \Gamma_{N,i}(k)=
\left[\begin{smallmatrix}
I \quad 0 \quad \ldots \quad 0 \quad 
\underbrace{(-1)^{k+1}I}_{(i+1)\text{-th block}} \quad
0 \quad \ldots \quad 0 \quad
\gamma_{Nk}^{0}I \quad
\ldots \quad
\gamma_{Nk}^{N-1}I \quad
0 \quad \ldots \quad 0
\end{smallmatrix}\right],$\\
for all $i=1,\cdots,q.$
\vspace{2mm}\\
In the following theorem we present the main result on a nonlinear matrix inequality which is essentially based on a~general result to derive a~storage functions for nonlinear systems~\cite{Sch17}. 

\begin{theorem}\rm
\label{thm:delay_dissipative}
Consider the time-delay system~\eqref{eq:cnt_time_delay_system_nonlinear} for some  $\tau_q\geq \ldots \geq \tau_1 > 0$ with respect to the quadratic supply rate $s$ given by \eqref{eq:supply_rate}. If there exists a function $\mathcal{S}_{N}$ given by \eqref{eq:storage_fixed_delay_nonlinear} such that for all $x_0,\ldots,x_q\in\mathbb{R}^n$ and $\xi^1_N,\ldots,\xi^q_N\in \mathbb{R}^{n(N+2)}$ 
\begin{align}
    \label{eq:nonlinear_LMI1}
    \begin{bmatrix}
    \Psi_{xx}(x_0,\ldots,x_q,\xi^1_N,\ldots,\xi^q_N) & \Psi_{xu}(x_0,\ldots,x_q)\\ \Psi_{xu}(x_0,\ldots,x_q)^\top & \Psi_{uu}(x_0)
\end{bmatrix}\leq 0,
\end{align}
then \eqref{eq:cnt_time_delay_system_nonlinear} is dissipative.
\end{theorem}
\begin{proof}
To verify the dissipativity, we consider the storage function \eqref{eq:storage_fixed_delay_nonlinear}, and show that 
\begin{align}
    \tfrac{\rm d}{ {\rm d}t}\mathcal{S}_{N}(x_t)&=\nabla\mathcal{S}_1(x(t))^\top\dot x(t) +
    2 \hat{x}_{N}(t)^\top P_{N} \dot{\hat{x}}_{N}(t)\nonumber\\ 
    &~~~~+\sum_{i=1}^q \tfrac{\rm d}{{\rm d}t} \int_{t-\tau_i}^t\Theta_i(x(s))\mathrm{d}s \label{eq:Theta_term} \\ &~~~~+ \sum_{i=1}^q  \tau_i \tfrac{\rm d}{{\rm d}t} \int_{t-\tau_i}^t\int_\theta^t \dot x(s)^\top R_i\dot{x}(s) \mathrm{d}s\mathrm{d}\theta  \label{eq:R_term} \\ 
&\leq s(y(t),u(t))=\begin{bmatrix}
    y(t)\\ u(t)
\end{bmatrix}^\top\begin{bmatrix}
    S_{11} & S_{12}\\ S_{12}^\top & S_{22}
\end{bmatrix}\begin{bmatrix}
    y(t)\\ u(t)
\end{bmatrix} \label{eq:dissipativity_estimate}
\end{align}
which should hold for all admissible inputs $u$ and all $t\geq 0$. 
Using the integration by parts and Lemma \ref{Lemma:Legendre_polynomials_properties}, we get
\begin{align*}
    \tfrac{\rm d}{{\rm d}t}\Omega^i_k(t) &= \int_{-\tau_i}^0 L^i_k(s) 
    \dot x_t(s) \, \mathrm{d}s\\
    &= \Big[ L^i_k(s) x(t+s) \Big]_{-\tau_i}^0 - \int_{-
    \tau_i}^0 
    \dot{L}^i_k(s) x(t+s) \, \mathrm{d}s \\
    &= L^i_k(0)x(t) - L^i_k(-\tau_i)x(t-\tau_i) - \int_{-\tau_i}^0 
    \dot{L}^i_k(s) x(t+s) \, \mathrm{d}s\\
    &=x(t) - (-1)^k x(t-\tau_i) - \scalebox{0.86}{$ \int_{-\tau_i}^0 \left( 
   \displaystyle
\sum_{j=0}^{k-1} \frac{2j+1}{\tau_i} (1 - (-1)^{k+j}) 
    L^i_j(s) \right) x(t+s) \, \mathrm{d}s$}\\
    &=x(t) - (-1)^k x(t-\tau_i) - \scalebox{0.86}{$ \displaystyle
\sum_{j=0}^{k-1} 
    \frac{2j+1}
    {\tau_i} (1 - (-1)^{k+j}) \underbrace{\int_{-\tau_i}^0 
    L^i_j(u) 
    x(t+s) \, \mathrm{d}s}_{\Omega^i_j(t)}$}\\
    &=x(t) - (-1)^k x(t-\tau_i) - \frac{1}{\tau_i} 
    \sum_{j=0}^{k-1} (2j+1) (1 - (-1)^{k+j}) \Omega^i_j(t)\\
    &=x(t) - (-1)^k x(t-\tau_i) + \frac{1}{\tau_i} 
    \sum_{j=0}^{k-1} \gamma^j_{Nk} \Omega^i_j(t).
\end{align*}
For $
\hat{x}_N(t) =G_N
\zeta(t)$ and $\dot{\hat{x}}_N(t)= H_N \zeta(t).$
Using the Leibniz rule for 
\eqref{eq:Theta_term} we obtain
\[
\tfrac{\rm d}{{\rm d}t}\sum_{i=1}^q\int_{t-\tau_i}^t\Theta_i(x(s))\mathrm{d}s=\sum_{i=1}^q\Theta_i(x(t))- \sum_{i=1}^q\Theta_i(x(t-\tau_i)).
\]

To compute the derivative of the term \eqref{eq:R_term}, we interchange the order of integration first 
\begin{align*}
\sum_{i=1}^q \tau_i \int_{t-\tau_i}^t\int_s^t\dot x(v)^\top R_i \dot x(v)\mathrm{d}v\mathrm{d}s&=\sum_{i=1}^q \tau_i \int_{t-\tau_i}^t\int_{t-\tau_i}^v\dot x(v)^\top R_i\dot x(v)\mathrm{d}s\mathrm{d}v\\&=\sum_{i=1}^q \tau_i \int_{t-\tau_i}^t(v-(t-\tau_i))\dot x(v)^\top R_i\dot x(v)\mathrm{d}v
\end{align*}
which allows for application of the Leibniz rule
{\footnotesize \begin{align}
\sum_{i=1}^q \tau_i\frac{\rm d}{{\rm d}t}\int_{t-\tau_i}^t\int_s^t\dot x(v)^\top R_i \dot x(v)\mathrm{d}v\mathrm{d}s&=\sum_{i=1}^q \tau_i \frac{\rm d}{{\rm d}t}\int_{t-\tau_i}^t(v-(t-\tau_i))\dot x(v)^\top R_i \dot x(v)\mathrm{d}v \nonumber \\
&=\footnotesize\sum_{i=1}^q \tau_i^2 \dot x(t)^\top R_i \dot x(t)-\sum_{i=1}^q \tau_i\int_{t-\tau_i}^t\dot x(s)^\top R_i \dot x(s)\mathrm{d}s. \label{eq:after_leibniz}
\end{align}}
The integral in \eqref{eq:after_leibniz} can be estimated using the Bessel--Legendre inequality \eqref{eq:bessel} from Lemma \ref{lemma:legendre}
\begin{align*}\nonumber
\sum_{i=1}^q \tau_i \int_{-\tau_i}^{0} \dot{x}(u)^\top R_i\dot{x}
(u)du &\geq \sum_{i=1}^q \xi_N^{i}(t)^\top\left[ 
\sum_{k=0}^{N} (2k+1) \Gamma_N(k)^\top R_i
\Gamma_N(k) \right] \xi^i_N(t)\\
&=\zeta(t)^\top \mathbf{\Gamma_N}^\top \mathbf{R_N}\mathbf{\Gamma_N}\zeta(t).
\end{align*}
This leads to 
{\footnotesize\begin{align}\nonumber
&~~~~   \sum_{i=1}^q \tau_i \frac{\rm d}{{\rm d}t}\int_{t-\tau_i}^t\int_s^t\dot x(v)^\top R_i\dot x(v)\mathrm{d}v\mathrm{d}s  \leq\sum_{i=1}^q \tau_i^2\dot{x}(t)^\top R_i \dot{x}(t)- \zeta(t)^\top\mathbf{\Gamma_N}^\top \mathbf{R_N}\mathbf{\Gamma_N}\zeta(t).\\
\nonumber
\end{align}}
We define 
\begin{align*}
\Psi(x_0,x_1,\ldots,x_q,\xi^1_N,\ldots,\xi^q_N)&:=\nabla\mathcal{S}_1(x_0)^\top(f(x_0,x_1,\ldots,x_q)+g(x_0)u)\\&+\zeta^\top(2\mathrm{sym}(G^\top_N P_N H_N)-\mathbf{\Gamma_N}^\top \mathbf{R_N}\mathbf{\Gamma_N})\zeta\\&+\sum_{i=1}^q \tau_i^2(f(x_0,x_1,\ldots,x_q)\\
&+g(x_0)u)^\top R_i (f(x_0,x_1,\ldots,x_q)+g(x_0)u)\\ &+\sum_{i=1}^q \Theta_i(x_0)- \sum_{i=1}^q\Theta_i(x_i) \\ 
&-h(x_0)^\top S_{11}h(x_0)-h(x_0)^\top S_{12}u-u^\top S_{12}^\top h(x_0)\\
&-u^\top S_{22}u.
\end{align*}
Then summarizing the above terms and using the assumption~\eqref{eq:nonlinear_LMI1}, we find that 
\footnotesize\begin{align}
\Psi(x_0,x_1,\ldots,x_q,\xi^1_N,\ldots,\xi^q_N,u)=\begin{bmatrix}
    1\\ u
\end{bmatrix}^\top\left[ \begin{smallmatrix}
    \Psi_{xx}(x_0,\ldots,x_q,\xi^1_N,\ldots,\xi^q_N) & \Psi_{xu}(x_0,\ldots,x_q)\\ \Psi_{xu}(x_0,\ldots,x_q)^\top & \Psi_{uu}(x_0)
\end{smallmatrix}\right]\begin{bmatrix}
    1\\ u
\end{bmatrix}\leq 0, \label{eq:Psi_estimate}
\end{align}
where $x_0=x(t),$ $x_i=x(t-\tau_i),$ where $i \in \{1,\ldots,q\}.$ 
Hence, \eqref{eq:Psi_estimate} implies \eqref{eq:dissipativity_estimate}. Integration of~\eqref{eq:dissipativity_estimate} over $[t_0,t_1]$ implies the dissipation inequality~\eqref{eq:dissip_ineq_cont_time}.
\end{proof}
If we take $P_N=R_i=0,$ $\forall i  \in \{1,\ldots,q\}$ in \eqref{eq:storage_fixed_delay_nonlinear}, we obtain the following storage function 
\begin{equation}\label{eq:storage_function_different_delays}
\begin{split}
\mathcal{S}(x_t) &= \mathcal{S}_1(x(t))+\sum_{i=1}^q 
\int_{t-\tau_i}^{t}  \Theta_i( x(s)) \mathrm{d}s,
\end{split}
\end{equation}
for some continuously differentiable $\mathcal{S}_1:\mathbb{R}^n\rightarrow[0,\infty)$ and continuous positive semi-definite function 
$\Theta_i:\mathbb{R}^n\rightarrow[0,\infty),$ which leads to a delay-independent dissipativity condition. 
\begin{corollary}
The system~\eqref{eq:cnt_time_delay_system_nonlinear} is dissipative with respect to the quadratic supply rate $s$ given by \eqref{eq:supply_rate}. If there exists a function $\mathcal{S}$ given by \eqref{eq:storage_function_different_delays} such that for all $x_0,\ldots,x_q\in\mathbb{R}^n$
\begin{align*}
    \begin{bmatrix}
    \Psi_{xx}(x_0,x_1,\ldots,x_q) & \Psi_{xu}(x_0)\\ \Psi_{xu}(x_0)^\top & \Psi_{uu}(x_0)
\end{bmatrix}\leq 0,
\end{align*}
where,
\begin{align*}
 \Psi_{xx}(x_0,x_1,\ldots,x_q)&:=\nabla\mathcal{S}_1(x_0)^\top f(x_0,x_1,\ldots,x_q)+\sum_{i=1}^q\Theta_i(x_0)-\sum_{i=1}^q\Theta_i(x_i)\\
 &-h(x_0)^\top S_{11}h(x_0),\\
  \Psi_{xu}(x_0)&:=\tfrac{1}{2}\nabla\mathcal{S}_1(x_0)^\top g(x_0)-h(x_0)^\top S_{12},\\
  \Psi_{uu}(x_0)&:=-S_{22}.
\end{align*}
\end{corollary}
\begin{proof}
 The result follows from the proof of Theorem \ref{thm:delay_dissipative}.
\end{proof}

\subsection{Interconnection of passive time-delay systems}
In the special case of impedance passive systems, the passivity is preserved under a particular type of negative feedback interconnections. For a similar result for nonlinear delay-systems using a slightly weaker notion of passivity we refer to \cite{kawano2023passivity}. 

We consider two impedance passive time-delay systems of the form
\begin{align}\label{eq:Interconnection}
\nonumber 
\dot x_j(t) &= f_j(x_j(t),x_j(t-\tau_{1,j}),\ldots, x_j(t-\tau_{q,j})) 
  + g_j(x_j(t))u_j(t), \quad t\geq 0,\\
y_j(t) &= h_j(x_j(t)),\quad \quad j=1,2,
\end{align}
with storage functions $\mathcal{S}_1$ and $\mathcal{S}_2$, respectively.
\begin{proposition}
If the systems \eqref{eq:Interconnection} are impedance passive  with storage functions $\mathcal{S}_1$ and $\mathcal{S}_2$ then the negative feedback interconnection
\[
u_1=v-y_2, \quad u_2=y_1,
\]
with external input $v$ is impedance passive with storage function $\mathcal{S}=\mathcal{S}_1+\mathcal{S}_2$.  
\end{proposition}
\begin{proof}
The impedance passivity inequalities hold individually for arbitrary inputs which gives
\begin{align*}
\mathcal{S}(x_{t_1})-\mathcal{S}(x_{t_0})&=\mathcal{S}_1(x_{t_1}^1)-\mathcal{S}_1(x_{t_0}^1)+\mathcal{S}_2(x_{t_1}^2)-\mathcal{S}_2(x_{t_0}^2) \\ 
&\leq \int_{t_0}^{t_1}y_1(s)^\top u_1(s)\mathrm{d}s+\int_{t_0}^{t_1}y_2(s)^\top u_2(s)\mathrm{d}s\\
&=\int_{t_0}^{t_1}y_1(s)^\top (v-y_2)(s)\mathrm{d}s+\int_{t_0}^{t_1}y_2(s)^\top y_1(s)\mathrm{d}s\\
&=\int_{t_0}^{t_1}y_1(s)^\top v(s)\mathrm{d}s.
\end{align*}
Hence, the interconnected system is impedance passive with respect to the input $v$ and the output $y_1$.
\end{proof}




\section{Dissipativity for linear time-delay systems}
\label{sec:linear}
In this section, we study the passivity of linear time-delay systems that are given by 
\begin{align}
\nonumber 
\dot x(t) &= Ax(t) + \sum_{i=1}^q A_ix(t-\tau_i)+ 
Bu(t),\quad  t\geq 0,\\
x(t) &= \phi(t),\quad  t \in [-\tau_q , 0], \label{eq:cnt_time_delay_system}
\\
y(t) &= Cx(t),\quad t \geq  0,\nonumber 
    \end{align}
where  
$A,A_i \in \mathbb{R}^{n \times n},$  where $i  \in \{1,\ldots,q\},$ and  $B \in \mathbb{R}^{n \times m},$ $C \in \mathbb{R}^{m \times n}$ and  $\tau_q \geqslant\ldots \geqslant \tau_1 > 0.$

Using Theorem~\ref{thm:delay_dissipative}, we present a sufficient condition for dissipativity for linear time-delay systems. We consider for some positive semi-definite $\Theta_i,R_i,\tilde P_N$ quadratic storage functions of the type 
\begin{align}\label{eq:storage_fixed_delay_linear}
\mathcal{S}_{N}(x_t)&=\tilde{x}_N(t)^\top \tilde P_N \tilde{x}_N(t)+\sum_{i=1}^q \int_{t-\tau_i}^t x(s)^\top \Theta_i x(s)\mathrm{d}s\\
&+ \sum_{i=1}^q \tau_i \int_{t-\tau_i}^t\int_s^t\dot x(v)^\top R_i \dot x(v)\mathrm{d}v\mathrm{d}s,\\\nonumber
\end{align} 
with the following extra-states 
{\footnotesize \begin{equation*}
\tilde {x}_{N}(t) = \begin{bmatrix} 
x(t)^\top &
\Omega_0^{1}(t)^\top  & 
\ldots &\Omega_{N-1}^{1}(t)^\top &
\ldots&
\Omega_{0}^{q}(t)^\top & \ldots &  \Omega_{N-1}^{q}(t)^\top 
\end{bmatrix}^\top , \hspace{2mm} \text{if} \hspace{2mm} N \geq 1
\end{equation*} }
and $\tilde{x}_0(t)=x(t)$ if $N=0.$
We consider 
\vspace{4mm}\\
$\tilde{G}_{N}=\begin{bmatrix}
I_n & 0&0&0 \\
   0& 0  & \rm diag(\tau_i I_{nN})_{i=1}^q & 0 \end{bmatrix}$ and
$\tilde{H}_{N}=\begin{bmatrix}
       \begin{bmatrix} A  & A_1 &  \ldots&  A_q &      0 & B \end{bmatrix} \\
       \tilde \Gamma_{N,1}(0) \\ \vdots \\ \tilde \Gamma_{N,q}(N-1) \end{bmatrix}. $
       \vspace{4mm}\\
Furthermore, we consider the following matrices
\begin{align*}
 M_0 &=
\begin{bmatrix}
       \sum_{i=1}^q \Theta_i - C^\top S_{11}C & 0 &0&-C^\top S_{12}\\
       0& - \rm diag(\Theta_i)_{i=1}^q & 0 &0  \\ 0&0&0&0\\
       -S_{12}^\top C&0&0&-S_{22}
\end{bmatrix},\\ 
M_{1,\tau}&=\left[\begin{smallmatrix}
  A^\top \sum_{i=1}^q \tau_i^2 R_i A & M_2^{12}& 0  &  A^\top \sum_{i=1}^q \tau_i^2 R_i  B \\
 *&  \begin{bmatrix}
A_i^\top  \sum_{i=1}^q \tau_i^2 R_i  A_j  
 \end{bmatrix}_{i,j=1}^q& 0 &(M_2^{42})^\top \\
 0&  0&0& 0\\
 *& 
*&0 &B^\top   \sum_{i=1}^q \tau_i^2 R_i  B 
\end{smallmatrix}\right]
\end{align*}
with
\begin{align*}M_2^{12}&= \begin{bmatrix}
 A^\top \sum_{i=1}^q \tau_i^2 R_i  A_1 &  \ldots &  A^\top \sum_{i=1}^q \tau_i^2 R_i  A_q 
\end{bmatrix}, \\ M_2^{42}&=\begin{bmatrix}
B^\top   \sum_{i=1}^q \tau_i^2 R_i  A_1 &  \ldots & B^\top  \sum_{i=1}^q \tau_i^2 R_i A_q 
 \end{bmatrix}.
\end{align*}
\begin{corollary}\rm
The linear time-delay system~\eqref{eq:cnt_time_delay_system} is dissipative for fixed $\tau_q\geq \ldots \geq \tau_1 > 0$ with respect to the quadratic supply rate $s$ given by \eqref{eq:supply_rate} if there exist positive semi-definite $\Theta_i,R_i\in\mathbb{R}^{n\times n}$ and positive semi-definite $\tilde P_N\in\mathbb{R}^{n(qN+1)\times n(qN+1)}$ which fulfill the following LMI
\begin{align}
\label{eq:LMI_fixed_tau}    
2 {\sym}( \tilde G_{N}^\top \tilde P_N \tilde H_{N})+ M_0 +M_{1,\tau}-\mathbf{\Gamma_N}^\top \mathbf{R_N}\mathbf{\Gamma_N}\leq 0.
\end{align}
\end{corollary}
\begin{proof}
To establish the 
dissipativity of system \eqref{eq:cnt_time_delay_system}, 
we consider the storage 
function \eqref{eq:storage_fixed_delay_linear} and show that 
the following inequality 
\begin{align}
    \tfrac{\rm d}{ {\rm d}t}\mathcal{S}_{N}(x_t)  
\leq s(y(t),u(t))=\begin{bmatrix}
    y(t)\\ u(t)
\end{bmatrix}^\top\begin{bmatrix}
    S_{11} & S_{12}\\ S_{12}^\top & S_{22}
\end{bmatrix}\begin{bmatrix}
    y(t)\\ u(t)
\end{bmatrix} 
\end{align}
holds for all admissible inputs $u$ and all $t\geq 0$. 
Let \begin{align*}
 \tilde\zeta(t)=\left[\begin{smallmatrix}x_t(0)^\top \quad x_t(-\tau_1)^\top \quad \ldots \quad x_t(-\tau_q)^\top\quad \left[\frac{1}{\tau_i}\Omega_0^{i}(t)^\top\quad \ldots \quad \frac{1}{\tau_i}\Omega_{N-1}^{i}(t)^\top\right]_{i=1}^q \quad u(t)
\end{smallmatrix}\right]^\top.\end{align*}
We have
\begin{align*}
    \tfrac{\rm d}{ {\rm d}t}\mathcal{S}_{N}(x_t)  
&=2 \tilde \zeta(t)^\top\sym( \tilde G_{N}^\top \tilde P_N\tilde H_{N}) \tilde\zeta(t)\\
&\footnotesize +\sum_{i=1}^q \tau_i^2 
\tilde \zeta(t)^\top\left[ \begin{smallmatrix}
    A & A_1 & \ldots & A_q &0&\ldots &B
\end{smallmatrix}\right]^\top R_i \left[\begin{smallmatrix}
    A & A_1 & \ldots & A_q &0&\ldots &B
\end{smallmatrix}\right]\tilde\zeta(t)
\\ 
&+\footnotesize\sum_{i=1}^q
x(t)^\top\Theta_i x(t)- \sum_{i=1}^q x(t-\tau_i)^\top\Theta_i x(t-\tau_i)\\
&-\footnotesize\sum_{i=1}^q \xi_N^{i \top}(t)  \left[ 
\sum_{k=0}^{N} (2k+1) \Gamma_N(k)^\top R_i 
\Gamma_N(k) \right] \xi^i_N(t)\\
&-2x(t)^\top C^\top S_{12}u(t)-u(t)^\top S_{22} u(t)-x(t)^\top C^\top S_{11} Cx(t)\\
&+2x(t)^\top C^\top S_{12}u(t)+u(t)^\top S_{22} u(t)+x(t)^\top C^\top S_{11} Cx(t)\\
&=\tilde \zeta(t)^\top(2 \sym( \tilde G_{N}^\top \tilde P_N \tilde H_{N})+ M_0 +M_{1,\tau}-\mathbf{\Gamma_N}^\top \mathbf{R_N}\mathbf{\Gamma_N})\tilde\zeta(t)\\
&+2x(t)^\top C^\top S_{12}u(t)+u(t)^\top S_{22} u(t)+x(t)^\top C^\top S_{11} Cx(t)\\
&\leq 2x(t)^\top C^\top S_{12}u(t)+u(t)^\top S_{22} u(t)+x(t)^\top C^\top S_{11} Cx(t).
\end{align*}
Using \eqref{eq:LMI_fixed_tau} and integrating \eqref{eq:dissipativity_estimate} over the interval $[t_0,t_1]$ yields the dissipation inequality~\eqref{eq:dissip_ineq_cont_time}.
\end{proof}
\begin{remark}
Note that in the impedance passive case we have $S_{22}=0$ which implies that $B^\top R_i B=0, \hspace{2mm} \forall i \in \{ 1,\ldots, q \}$ must hold. This means that $\ran B\subseteq \bigcap\limits_{i=1}^{q} \ker R_i$ which poses some restrictions on the choice of $R_i$. In particular, under the very reasonable assumption that $B\neq 0$ the matrices $R_i$ cannot be positive definite.
\end{remark}
\section{Conclusion}
\label{sec:conclusion}
In this paper, we studied dissipativity properties for a class of nonlinear time-delay systems with multiple constant delays. By employing Lyapunov--Krasovskii storage functionals together with Bessel--Legendre integral inequalities, we derived sufficient delay-dependent dissipativity conditions with respect to general quadratic supply rates. A delay-independent criterion was also obtained as a special case by considering a simplified storage functional.

For linear time-delay systems, the proposed framework led to tractable sufficient conditions expressed in terms of linear matrix inequalities. These conditions can be checked numerically using standard semidefinite programming tools and provide a systematic way to verify dissipativity, including passivity-type properties, in the presence of multiple delays.

Future work will focus on reducing the conservatism of the obtained conditions and on developing efficient numerical methods for the nonlinear case, where the resulting dissipativity conditions are generally nonconvex. Another relevant direction is the application of the proposed criteria to concrete physical and engineering systems with delayed feedback or delayed damping.

\section*{Acknowledgments}
This work was funded by the Deutsche Forschungsgemeinschaft (DFG, German
Research Foundation) – Project-ID 531152215 – CRC~1701. Ikram El Haskouki gratefully acknowledges University of Wuppertal, Germany, for hosting her during her research stay and for their support.

\bibliographystyle{abbrv}
\bibliography{references_P.bib}

\end{document}
